\par\Needspace{6\baselineskip}\medskip
\hypertarget{native-141-article101}{}\pdfbookmark[2]{101.}{article-101}
\begin{center}101.\end{center}
Introducendo quantitatem $\zeta$ statuendoque $\BareRoot{(l-z)}=\dfrac my$ vel $\BareRoot{(L+z)}=\dfrac MY$, nec non
\[
\mathrm{15)}\qquad\frac{mm}{\tfrac56+l+\zeta}=h
\]
\noindent vel
\[
\mathrm{15^*)}\qquad\frac{MM}{L-\tfrac56-\zeta}=H,
\]
\noindent aequationes 13, $13^*$ hancce formam induunt
\[
\begin{aligned}
\mathrm{16)}\qquad&\frac{(y-1)yy}{y+\tfrac19}=h\\[.6em]
\mathrm{16^*)}\qquad&\frac{(Y+1)YY}{Y-\tfrac19}=H,
\end{aligned}
\]
\noindent adeoque omnino identicae fiunt cum iis ad quas in ellipsi perventum est (15, $15^*$ art. 91). Hinc igitur, quatenus $h$ vel $H$ pro cognita haberi potest, $y$ vel $Y$ deduci poterit, ac dein erit
\[
\begin{aligned}
\mathrm{17)}\qquad&z=l-\frac{mm}{yy}\\[.6em]
\mathrm{17^*)}\qquad&z=\frac{MM}{YY}-L.
\end{aligned}
\]
\noindent Ex his colligitur, omnes operationes supra pro ellipsi praescriptas perinde etiam pro hyperbola valere, donec e valore approximato ipsius $h$ vel $H$ eruta fuerit quantitas $y$ vel $Y$; dein vero quantitas $\dfrac{mm}{yy}-l$ vel $L-\dfrac{MM}{YY}$, quae in ellipsi positiva evadere debebat, in parabolaque $=0$, fieri debet negativa in hyperbola: hoc itaque criterio genus sectionis conicae definietur. Ex inventa $z$ tabula nostra dabit $\zeta$, hinc orietur valor correctus ipsius $h$ vel $H$, cum quo calculus repetendus est, donec omnia ex asse conspirent.
\par
Postquam valor verus ipsius $z$ inventus est, $c$ inde per formulam $c=1+2z+2\BareRoot{(z+zz)}$ derivari posset, sed praestat, etiam ad usus sequentes, angulum auxiliarem $n$ introducere, per aequationem $\operatorname{tang}2n=2\BareRoot{(z+zz)}$ determinandum; hinc fiet $c=\operatorname{tang}2n+\BareRoot{(1+\operatorname{tang}2n^2)}=\operatorname{tang}(45^0+n)$.
